\newpage
\begin{center}
	{ВВЕДЕНИЕ}
\end{center}

Современные спортивные организации, осуществляющие подготовку спортсменов по различным видам спорта, сталкиваются с необходимостью обработки значительных объемов информации. Это и учет спортсменов и тренеров, и составление расписания тренировок, и регистрация соревнований и их результатов, и присвоение спортивных разрядов, и контроль движения инвентаря. Автоматизация всего этого облегчила и систематизировала бы процесс.

Традиционные методы ведения документации, основанные на использовании бумажных журналов и редких электронных таблиц, обладают рядом существенных недостатков. Дублирование информации, высокая трудоемкость формирования отчетности, отсутствие оперативного контроля за движением инвентаря и невозможность одновременной работы нескольких пользователей снижают эффективность управления спортивной организацией. Особенно остро эти проблемы проявляются в небольших и средних учреждениях, не располагающих бюджетами на приобретение дорогостоящих специализированных решений.

Анализ существующих на рынке программных продуктов выявил их ограниченную применимость для нужд небольшой спортивной школы. Специализированные решения, такие как «АИС Мой спорт» и «SportBase», ориентированы на государственную отчетность и крупные федерации, имеют высокую стоимость внедрения и сопровождения либо не поддерживают детальный учет спортивного инвентаря. Универсальные учетные системы (например, «1С:Бухгалтерия» или «Парус») не учитывают спортивную специфику: в них отсутствуют классификаторы видов спорта и разрядов, механизмы составления расписания с привязкой к тренерам. Использование электронных таблиц Microsoft Excel, хотя и не требует финансовых затрат, приводит к дублированию информации и невозможности одновременной работы нескольких пользователей. Таким образом, существующие решения не обеспечивают комплексной автоматизации спортивной школы при приемлемой стоимости.

Поэтому разработка собственной информационной системы, адаптированной под нужды конкретной спортивной организации и реализованной на доступной платформе, является актуальной задачей.

\textit{Цель настоящей работы} - разработка информационной системы на платформе 1С:Предприятие 8.3 для комплексного учета деятельности спортивной организации. Для достижения поставленной цели необходимо решить \textit{следующие задачи:}

\begin{itemize}
	\item провести анализ деятельности спортивной организации и обосновать необходимость автоматизации ключевых бизнес-процессов;
	\item спроектировать структуру информационной системы, включающую справочники, документы, регистры и отчеты;
	\item разработать и реализовать документы, обеспечивающие планирование тренировочного процесса, организацию и проведение соревнований, оформление спортивных разрядов, а также складской учет спортивного инвентаря;
	\item реализовать алгоритм заполнения планировщика расписания и функцию импорта данных из файлов Excel;
	\item реализовать разграничение прав доступа на основе ролей «Администратор» и «Тренер» и провести функциональное тестирование системы.
\end{itemize}

\textit{Предметом исследования} являются структура объектов конфигурации 1С (справочников, документов, регистров), алгоритмы программного заполнения планировщика расписания и импорта данных из внешних файлов, а также механизмы разграничения прав доступа пользователей.

\textit{Объектом разработки} является информационная система на платформе 1С:Предприятие 8.3 для комплексного учета деятельности спортивной организации.


\textit{Структура и объем работы.} Выпускная квалификационная работа состоит из введения, 4 разделов основной части, заключения, списка использованных источников, 2 приложений. Текст выпускной квалификационной работы равен \formbytotal{lastpage}{страниц}{е}{ам}{ам}.

\textit{Во введении} сформулирована цель работы, поставлены задачи разработки, описана структура работы, приведено краткое содержание каждого из разделов.

\textit{В первом разделе} на стадии анализа предметной области рассматривается деятельность спортивной организации, выявляются ключевые бизнес-процессы, подлежащие автоматизации, анализируются существующие программные решения и обосновывается необходимость разработки собственной информационной системы.

\textit{Во втором разделе} на стадии технического задания формулируются основание для разработки, цель и назначение системы, требования к функциональности, интерфейсу, надежности, информационной безопасности и программной совместимости, а также моделируются варианты использования для ролей «Администратор» и «Тренер».

\textit{В третьем разделе} на стадии технического проектирования представлены общая архитектура системы, обоснование выбора платформы, структура конфигурации, описание справочников, документов и регистров, а также алгоритмы программной реализации.

\textit{В четвертом разделе} приводится спецификация программных модулей, описание реализации ключевых алгоритмов, результаты функционального тестирования и инструкция по развертыванию системы.

\textit{В заключении} излагаются основные результаты работы, полученные в ходе разработки.

\textit{В приложении А} представлен графический материал. \textit{В приложении Б} представлены фрагменты исходного кода программы.